\section{The original finite source and the degree-zero DI summand}
\label{sec:di-finite-source-assembly}

Throughout, $f:X\to\operatorname{Spec}k$ is smooth of pure relative
dimension $n$, $k$ has characteristic $p$, $Y=X^{(p)}$, and
$K=F_{X/k*}\Omega^\bullet_{X/k}$ is the original natural-degree complex.
The Cartier isomorphism used here is the actual global isomorphism already
constructed on these sheaves. No new choice of homology object is made.

\begin{definition}[The original degree-zero summand]
\label{def:di-actual-zero-summand}
\lean{EnvelopeProperty.DI.FiniteSourceAssembly.zeroComponent,EnvelopeProperty.DI.FiniteSourceAssembly.zeroComponent_closed,EnvelopeProperty.DI.FiniteSourceAssembly.zeroCochainMap,EnvelopeProperty.DI.FiniteSourceAssembly.zeroCochainMap_f_zero,EnvelopeProperty.DI.FiniteSourceAssembly.zeroTotalMap}
\leanok
\uses{thm:cartier-geometric-global-sheaf-isomorphism}
There are no incoming boundaries in degree zero of $K$. Thus the canonical
quotient $\pi:Z^0K\to H^0K$ is an isomorphism. Define
\[
 a_0:\Omega^0_{Y/k}\xrightarrow{C^{-1}}H^0K
       \xrightarrow{\pi^{-1}}Z^0K\hookrightarrow K^0.
\]
This gives a cochain map from the single complex in degree zero to $K$.
For any family of opens $\mathcal U$, compose with the original total
\v Cech augmentation to obtain its degree-zero map into
$\operatorname{Tot}\check C(\mathcal U,K)$.
\end{definition}
\begin{proof}
The composite of the cycle inclusion with the outgoing differential is
zero, so $a_0$ is closed. The universal construction of maps from a single
complex gives the stated cochain map, whose degree-zero component is
exactly the original single-object identification followed by $a_0$.
Composition with the existing augmentation preserves the cochain equation.
\end{proof}

\begin{lemma}[The actual degree-zero Cartier class]
\label{lem:di-actual-zero-cartier-class}
\lean{EnvelopeProperty.DI.FiniteSourceAssembly.zeroCochainMap_class}
\leanok
\uses{def:di-actual-zero-summand}
The map induced by the constructed degree-zero cochain map on homology,
using the canonical homology identification of its original single
complex, is exactly the entire global inverse Cartier map in degree zero.
\end{lemma}
\begin{proof}
The map on cycles is the original single-cycle identification followed by
$C^{-1}\pi^{-1}$. Naturality of the original quotient from cycles to
homology therefore identifies its homology class with
$C^{-1}\pi^{-1}\pi=C^{-1}$. This is an equality of the entire sheaf maps,
before evaluation on any open.
\end{proof}

\begin{lemma}[The finite source is already below the cutoff]
\label{lem:di-finite-source-good-truncation}
\lean{EnvelopeProperty.DI.FiniteSourceAssembly.source_extend_isStrictlyLE,EnvelopeProperty.DI.FiniteSourceAssembly.sourceTruncationIso,EnvelopeProperty.DI.FiniteSourceAssembly.truncatedMap,EnvelopeProperty.DI.FiniteSourceAssembly.truncatedMap_inclusion}
\leanok
For any abelian category and original objects $(M_q)_{q<p}$, put
$S=\bigoplus_{q<p}M_q[-q]$, first as a natural-degree complex and then
extended by zero to integer degrees. The original good-truncation
inclusion $\tau_{<p}S\to S$ is an isomorphism. Any cochain map $a:S\to T$
has the canonical factorization
\[
 S\xrightarrow{\iota_S^{-1}}\tau_{<p}S
   \xrightarrow{\tau_{<p}a}\tau_{<p}T,
\]
whose composite with the original inclusion into $T$ is exactly $a$.
This assertion includes $p=0$ and assumes no invertibility or
quasi-isomorphism property of $a$.
\end{lemma}
\begin{proof}
Every summand of $S$ is zero in each degree $q\ge p$, so the existing
finite-source term computation makes the original coproduct zero there.
Extension by zero retains these terms and adds only zero negative terms.
Consequently $S$ is strictly supported in degrees at most $p-1$, and the
native good-truncation inclusion is invertible. Its naturality square for
$a$, composed with its inverse on the source, gives the entire asserted
factorization identity.
\end{proof}

These constructions supply the zero summand and the source truncation for
the final finite-source assembly. They do not assert that the positive
summands, their Cartier comparison, or the global DI decomposition have
already been constructed.
